\section{Efficient Computation}
\label{sec:computation}
Storing dense attention matrices requires space for all token pairs. \dt{} computes local transpose products in small blocks and accumulates their results directly into the input slopes. The attention and memory rules in Sections~\ref{sec:attention}--\ref{sec:memory} support tiled and chunked implementations.

\paragraph{Attention in tiles.}
A tile is a small block of query and source positions. Following FlashAttention \citep{dao2022flashattention}, \dt{} reconstructs the two executions' attention probabilities for one tile, uses them to update query, key, and value slopes, and reuses the workspace for the next tile. The softmax rule needs a weighted average over the whole row. One sweep collects those row statistics. Subsequent sweeps use them when propagating through each tile. For $N$ input-plus-response positions and head dimension $d$, the stored vectors and row statistics occupy $O(Nd)$ auxiliary space per head. Dense attention requires $O(N^2d)$ arithmetic.

\paragraph{Memory in chunks.}
A chunk is a consecutive group of tokens. The memory state at a chunk boundary summarizes what earlier tokens passed forward. In reverse, the slopes for that state pass back to the preceding chunk. Within each chunk, matrix products propagate the slopes through its reads, writes, and gates. This organization uses Flash Linear Attention (FLA) \citep{yang2024fla}. At fixed state dimensions and chunk size, the stored token slopes and boundary states grow linearly with sequence length.

The complete attribution call includes two forward executions and one reverse pass. It recomputes intermediate activations one layer at a time during the reverse pass, so each layer's workspace can be reused for the next. Appendix~\ref{app:computation} details the storage accounting and accelerated operations. Section~\ref{sec:evaluation} measures the complete call's time and memory.
